\scanpage{37}

pour $n$ variable. Ceci dit, on a un hom évident

\[
\mathcal E=(\unsure{\mathfrak S_3|\mathscr S_3}\times\Gamma)\underset{1/2}{\cdot}\Pi_\infty\longrightarrow
\widetilde{\mathcal E}=(\mathfrak S_3\times\mathcal A)\underset{1/2}{\cdot}\Pi_\infty
\tag*{(41)}
\]

déduit de (40), et les opérations de $\mathcal E$, \unsure{sur|en} les ens. précédents se font \uline{à travers} $\widetilde{\mathcal E}$ ....

(1.9) Les développements précédents n'utilisaient pas que \deleted{$\mathbb Q$}la clôture alg. \inserted{envisagée} $\overline{\mathbb Q}$ de $\overline{\mathbb Q}$ était plongée dans $\mathbb C$ -- ils marcheraient $=$ en remplaçant $\mathbb Q$ par n'importe quel corps premier, à condition de se borner à des $n$ premiers à la car. ($\overline C_\infty$ est alors le revêtement universel abélien \uline{modéré} de $(\overline C_1,I)$ ...). Il faut l'utiliser, pour faire le lien avec les cartes $3$-pondérées finies abéliennes, et le groupe cartographique.

On va d'abord \unsure{pointer|ponctuer|trivialiser} $\overline C_\infty$ au dessus de $\unsure{\sigma|\alpha|\varpi}=-\bar j\in\omega=\{\sigma,\bar\sigma\}$, pour pouvoir l'identifier au rev. universel profini \deleted{pointé} \uline{sur $\sigma$}.

% archive sheet number: 36
